\section{Modus Ponens dan Modus Tollens}

Di antara aturan inferensi, \textit{Modus Ponens} dan \textit{Modus Tollens} termasuk yang paling sering dipakai, baik dalam matematika maupun dalam penalaran sehari-hari.

\subsection{Modus Ponens}

Istilah \textit{modus ponendo ponens} (bahasa Latin) merujuk pada cara menetapkan konsekuen dari implikasi yang telah dianggap benar.

\begin{teorema}[Modus Ponens]
Jika $p \rightarrow q$ benar dan $p$ benar, maka $q$ benar.
\end{teorema}

\textbf{Skema argumen Modus Ponens:}
\[
\begin{array}{cl}
p \rightarrow q & \text{(Premis: implikasi)} \\
p & \text{(Premis: anteseden benar)} \\
\hline
\therefore q & \text{(Konklusi: konsekuen benar)}
\end{array}
\]

\begin{contoh}
\begin{itemize}
    \item \textbf{Premis 1 ($p \rightarrow q$)}: ``Jika Rini memperoleh skor TOEFL 500, maka Rini memenuhi persyaratan yudisium.''
    \item \textbf{Premis 2 ($p$)}: ``Rini memperoleh skor TOEFL 500.''
    \item \textbf{Konklusi ($\therefore q$)}: ``Rini memenuhi persyaratan yudisium.''
\end{itemize}
\end{contoh}

Ekspresi $[(p \rightarrow q) \land p] \rightarrow q$ merupakan tautologi, sehingga skema Modus Ponens selalu menghasilkan argumen yang valid.

\subsection{Modus Tollens}

Istilah \textit{modus tollendo tollens} merujuk pada penolakan anteseden ketika konsekuen dari implikasi ditetapkan salah. Aturan ini sejalan dengan ekuivalensi kontraposisi $p \rightarrow q \equiv \neg q \rightarrow \neg p$ yang telah dibahas pada bab sebelumnya \cite{rosen2019}.

\begin{teorema}[Modus Tollens]
Jika $p \rightarrow q$ benar dan $\neg q$ benar, maka $\neg p$ benar.
\end{teorema}

\textbf{Skema argumen Modus Tollens:}
\[
\begin{array}{cl}
p \rightarrow q & \text{(Premis: implikasi)} \\
\neg q & \text{(Premis: konsekuen salah)} \\
\hline
\therefore \neg p & \text{(Konklusi: anteseden salah)}
\end{array}
\]

\begin{contoh}
\begin{itemize}
    \item \textbf{Premis 1 ($p \rightarrow q$)}: ``Jika si penjaga kolam adalah pelaku, maka sepatu bot miliknya basah oleh lumpur.''
    \item \textbf{Premis 2 ($\neg q$)}: ``Sepatu bot milik penjaga kolam kering (tidak basah oleh lumpur).''
    \item \textbf{Konklusi ($\therefore \neg p$)}: ``Penjaga kolam bukan pelaku.''
\end{itemize}
\end{contoh}

Ekspresi $[(p \rightarrow q) \land \neg q] \rightarrow \neg p$ merupakan tautologi, sehingga skema Modus Tollens valid.
